\begin{abstract}
Physical properties are often generated, calculated, and interpreted in
different representations of a theory.  It does not follow that a property
made manifest in one description has been established for the physical
theory.  We develop a proof-bearing calculus for representation change in
which a morphism records, separately, exact preservation, quantitative
approximation, loss from the target variables, theorem-backed recovery, and
the index domain on which these claims are uniform.  The formal core is a
uniform recovery-to-spectrum theorem.  For positive semigroups, convergence
of Euclidean quadratic forms on a dense centered observable core, together
with one index-independent exponential rate, excludes target spectrum below
that rate.  Exact countermodels show that the conclusion fails as a physical
gap inference when the core is incomplete, finite-index rates collapse, or
cutoff and volume limits are interchanged without uniform recovery.  We use
four obstructions in lattice gauge reasoning to demonstrate that quantifier
domains, boundary charge, adaptive history, and partition-function
normalization are distinct forms of representational loss requiring different
repairs.  A complete Yang--Mills representation chain then identifies the
open interfaces between regulated local control, continuum Schwinger data,
Osterwalder--Schrader reconstruction, and a Hamiltonian mass gap.  The result
is a formal methodology for cross-representation warrant, compatible with
several realist and empiricist ontologies; it supplies no new construction of
four-dimensional continuum Yang--Mills theory.
\end{abstract}
